% !TEX root = ../../main.tex
\section{مناقشة النتائج} \label{part1:sec4}
% \thispagestyle{empty}

من خلال هذه الدراسة والمحاكاة العملية باستخدام برنامج \el{\textbf{MATLAB R2023b}}، خرجنا بالنتائج التالية:

\begin{enumerate}
    \item أثبتت مرشحات التردد الرنيني \el{(\textbf{Tuned Filters})} كفاءة عالية في امتصاص تيارات التوافقيات عند ترددات الرنين المحددة، مما أدى لخفض نسبة \el{$THD$} بشكل تدريجي مع كل إضافة.
    \item ساهمت المرشحات في تحسين معامل الاستطاعة وتقريب شكل موجة التيار والجهد من الشكل الجيبي النقي، مما يقلل من الفواقد والحرارة في عناصر الشبكة.
    \item بفضل إضافة مرشحات الرتب 7 و 11، تمكنا من الحفاظ على جودة الجهد ضمن الحدود العالمية المسموح بها (غالباً أقل من 3\% أو 5\% حسب المواصفة).
    \item أكدت الدراسة أن للحصول على نتائج صحيحة ومطابقة للواقع الفيزيائي يجب أن ضبط إعدادات المحاكاة بشكل دقيق و إدخال بارامترات العناصر بعناية شديدة، حيث أن أي خطأ في القيم أو الإعدادات قد يؤدي لنتائج غير دقيقة أو غير منطقية.
\end{enumerate}
